\documentclass[10pt]{article}

\usepackage[utf8]{inputenc}

\usepackage[T1]{fontenc}

\usepackage{amsmath}

\usepackage{amsfonts}

\usepackage{amssymb}

\usepackage[version=4]{mhchem}

\usepackage{stmaryrd}

\usepackage{graphicx}

\usepackage[export]{adjustbox}

\graphicspath{ {./images/} }



\begin{document}

ры кривизн многообразия $M$. Если $M-$ сфера, то единственная Ф. т. $N$ - центр сферы (фокус). Ф. т. определяется также в терминах якобиевых полей (см. [2]). Множество всех Ф. т. называется каус тикой $N$. Ф.т. и каустики играют важную роль в различных задачах физики (оптика, квантовая механика, электродинамика и т. д.).



Лит.: [1] М асло в В. П., Фед ор юк М. В., Квазиклассическое приближение для уравнений квантовой механики, М., 1976; [2] Б и шоп Р., К риттенде н Р., Геометрия многообразий, пер. с англ., M., 1967.



М. В. Федорюк. ФОККЕРА - ПЛАНКА УРАВНЕНИЕ - см. Диффузионный процесс, Колмогорова уравнение, Марковский процесс, Маркова цепь, Маркова цепь с непрерывным временем, Орнитейна - Уленбека процесс.



\begin{center}

\includegraphics[max width=\textwidth]{2023_07_08_bb557379dc09012f5b44g-1}

\end{center}



Фокус. ФOKУC - один из видов особых точек обыкновенных дифференциальных уравнений. Все интегральные кривые, проходящие через точки достаточно малой окрестности такой особой точки, представляют собой спирали с бесконечным числом витков, неограниченно приближающихся к особой точке, навиваясь на нее (рис.).



В. И. Битюцков. ФOЛЬKA - БPУXA HEPABEНСТВО - для любого оператора $\hat{\boldsymbol{A}}$ и любого самосопряженного оператора $\hat{\boldsymbol{H}}$ неравенство:



\[

b(\hat{A}) \geqslant g(\hat{A}) f(c(\hat{A}) / 4 g(\hat{A})),

\]



где



\[

g(\hat{A})=\frac{1}{2} \operatorname{Tr}\left(\left(\hat{A}^{+} \hat{A}+\hat{A} \hat{A}^{+}\right) e^{-\beta \hat{H}}\right),

\]



$b(\hat{A})=Z^{-1} \int_{0}^{1} \operatorname{Tr}\left(\hat{A}^{+} e^{-\beta x \hat{H}} \hat{A} e^{-\beta(1-x) \hat{H}}\right) d x, \quad Z=\operatorname{Tr} e^{-\beta \hat{H}}$.



\[

c(\hat{A})=\operatorname{Tr}\left(\left[\hat{A}^{+}[\beta \hat{H}, \hat{A}]\right] e^{-\beta \hat{H}}\right),

\]



a $f$ - функция, задаваемая соотношением



\[

f(x \operatorname{th} x)=x^{-1} \operatorname{th} x .

\]



Ф.- Б. н. нашло применение в проблеме фазовых переходов. Jum.: [1] Falk H., B ruch D. W., «Phys. Rev.», 1969, v. 180, № 1, p. 442; [2] Dys on F., Lieb E. H., S im on B., "Phys. Rev. Lett.", 1976, v. 37, № 3, p. 120 . Д. П. Санкович. ФOHOH - квазичастица с энергией $E=\hbar \omega$ и импульсом $\boldsymbol{p}=\hbar \boldsymbol{k}$, которая в соответствии с основными принципами квантовой механики сопоставляется волне смещений атомов (ионов, молекул и т. п.) материальной среды с частотой $\omega$ и волновым вектором $\boldsymbol{k}$ (здесь $\hbar-$ постоянная Планка). Является основным понятием квантовой теории колебаний материальных сред, к-рые могут рассматриваться в виде газа $\Phi$.



Лит.: [1] Т а м м И. Е., Собрание научных трудов, т. 1, М., 1975, c. 168-85; [2] Бор н М., Хуан Кунь, Динамическая теория кристаллических решеток, пер. с англ., М., 1958. Е.Д. Белоколос. ФOPMAЛЬНЫЙ ИНТЕГPАЛ анал итиче ской с ис темы дифференциальных уравнений- формальный ряд, производная которого в силу этой системы равна 0 (в смысле равенства формальных рядов).



A. И. Нeйumadm. ФОРмАЛЬНЫй РяД (степенной) от переменных $x_{1}, \ldots$, $x_{n}$ над полем $P$ - алгебраическое выражение вида



\[

F\left(x_{1}, \ldots, x_{n}\right)=\sum_{k=0}^{\infty} F_{k}\left(x_{1}, \ldots, x_{n}\right)

\]



где $F_{k}-$ форма от $x_{1}, \ldots, x_{n}$ степени $k$ с коэффициентами из $P$. Обычно $P$ - поле действительных или комплексных чисел. Коэффициенты могут быть функциями нек-рых параметров. При естественном определении операций сложения и умножения Ф. р. образуют коммутативное кольцо с единицей. По обычным правилам вводятся производные $\Phi$. р. по переменным $x_{i}$ или по параметрам. Можно определить также композицию Ф. р., считая $x_{i}$ Ф. р. от новых переменных $z_{j}$. Ф. р. широко используются в возмуцений теории.



Лит.: [1] Би ркгоф Дж. Д., Динамические системы, пер. с англ., М.-Л., 1941. ФOPMФAKTOP в теории элементарных част и ц - функция, описывающая влияние протяженности частицы на ее взаимодействие с другими частицами и полями. Термин «Ф.» заимствован из теории рассеяния рентгеновских лучей, а его применение основано на наглядном представлении о том, что, напр., протон проводит часть времени в виртуальном состоянии «нейтрон $+\pi^{-}-$мезон». Поэтому заряд его оказывается "размазанным» с нек-рой плотностью $е \rho(r)$. Тогда, напр., амплитуда рассеяния электронов на таком размазанном протоне отличается от амплитуды рассеяния на точечном протоне множителем, называемым фо р м фа к т о ро м протона



\[

F(q)=\int \rho(r) \exp (i q r) d V

\]



где $q$ - передаваемый при рассеянии импульс.



В последовательной релятивистской локальной теории реальное размазывание невозможно, а строгий смысл термина Ф. в ней определяется следуюшим рассуждением. Плотность энергии взаимодействия электромагнитного поля, описываемого 4-потенциалом $A_{\mu}(x)$ со свободным фермионом, волновая функция к-рого $\psi(x)$, имеет вид:



\[

H(x)=i e \overline{\psi(x)} \gamma^{\mu} \psi(x) A_{\mu}(x) \equiv j_{0}^{\mu}(x) A_{\mu}(x),

\]



где $\gamma^{\mu}-$ Дирака матрицы, а $j_{0}^{\mu}-$ функция, называемая электромагнитным током свободных фер м и о но в. Но само взаимодействие меняет оператор тока $j_{0}^{\mu}$. Матричный элемент электромагнитного тока взаимодействующего протона, взятый между состояниями реального протона с 4-импульсами $p$ и $p^{\prime}$, с учетом релятивистской инвариантности, Дирака уравнения и сохранения заряда, в общем случае можно записать в виде:



\[

\left\langle p^{\prime} \mid j^{\mu}(0) p\right\rangle=\left(4 p_{0} p_{0}^{\prime}\right)^{-1} \bar{u}\left(p^{\prime}\right)\left\{\gamma^{\mu} F_{1}\left(q^{2}\right)+i \sigma^{\mu \nu} q_{v} F_{2}\left(q^{2}\right)\right\} u(p),

\]



где $q=p-p^{\prime}, \sigma^{\mu \nu}=\frac{i}{2}\left(\gamma^{\mu} \gamma^{\nu}-\gamma^{\nu} \gamma^{\mu}\right)$. Входящие сюда скалярные функции $F_{1}\left(q^{2}\right)$ и $F_{2}\left(q^{2}\right)$ называются э ле к т р иче с ким и магнитным Ф. протона; о них заранее можно утверждать лишь то, что в пределе $q \rightarrow 0$ для длинных волн или малых передаваемых импульсов $F_{1}(0)=e$, где $e-$ наблюдаемый заряд, а $F_{2}(0)=\mu$, где $\mu-$ полный магнитный момент протона. Для свободной частицы $F_{1}\left(q^{2}\right)=e$, а $F_{2}\left(q^{2}\right)=\mu_{0}$, где $\mu_{0}-$ «нормальный» магнитный момент дираковской частицы с зарядом $e, \mu_{0}=\frac{e \hbar}{2 m c}$. В системе координат, где $q_{0}=0$, выражения



\[

\rho_{1,2}=(2 \pi)^{-2} \int d q \exp (\text { iqr }) F_{1,2}\left(-q^{2}\right)

\]



можно считать пространственными распределениями соответственно заряда и магнитного момента взаимодействующей частицы. Благодаря Ф. $F_{1}$ и $F_{2}$ взаимодействующий протон выглядит протяженным; однако нельзя говорить о реальном физич. размазывании протона, поскольку взаимодействующий ток остается локальным оператором и условие микропричинности не нарушается. Аналогично электромагнитным $\Phi$. $F_{1}$ и $F_{2}$ можно ввести мезонные $\Phi$.





\end{document}